\chapter{三角函数：单位圆}\label{ch:g11:trigo}

三角\omterm{def:g11:func:function}{函数}始于直角三角形，但其真正的归宿是
\emph{单位圆}，在其上\omterm{def:g11:trigo:cossin}{余弦}与\omterm{def:g11:trigo:cossin}{正弦}成为任意
\omterm{def:g10:numbers:sets}{实数}的\omterm{def:g11:func:function}{函数}。本章引入\omterm{def:g11:trigo:radian}{弧度}、\omterm{def:g10:numbers:sets}{实数}轴在圆周上的卷绕、
需要熟记的特殊值，以及生成全部经典恒等式的对称性。
关于 $\cos$ 与 $\sin$ 作为\omterm{def:g11:func:function}{函数}的研究——
即其变化与\omterm{def:g11:deriv:derivative}{导数}——
在\cref{ch:g12:trigo}中进行。

\section{弧度与卷绕}

\begin{definition}[弧度]\label{def:g11:trigo:radian}
设 $\mathcal C$ 为半径为 $1$、以原点 $O$ 为圆心的圆，
并设 $I(1, 0)$。圆心角的\emph{弧度}\index{弧度}量
是顶点位于 $\mathcal C$ 圆心处的角在 $\mathcal C$ 上所截弧的长度。
由于整个圆的周长为 $2\pi$：
\[
360^\circ = 2\pi \text{ 弧度},
\qquad
180^\circ = \pi \text{ 弧度},
\qquad
1 \text{ 弧度} = \frac{180^\circ}{\pi} \approx 57.3^\circ .
\]
\end{definition}

\begin{method}[换算]\label{met:g11:trigo:convert}
角度与\omterm{def:g11:trigo:radian}{弧度}成正比：由角度化为\omterm{def:g11:trigo:radian}{弧度}乘以 $\frac{\pi}{180}$，
由\omterm{def:g11:trigo:radian}{弧度}化为角度则乘以 $\frac{180}{\pi}$。例如
$60^\circ = 60 \times \frac{\pi}{180} = \frac{\pi}{3}$，且
$\frac{3\pi}{4} = \frac{3 \times 180^\circ}{4} = 135^\circ$。
\end{method}

\begin{definition}[实数轴在圆上的卷绕]\label{def:g11:trigo:winding}
对每个\omterm{def:g10:numbers:sets}{实数} $t$，对应点 $M(t)$，该点位于 $\mathcal C$ 上，
通过沿圆周走过长度 $\abs t$ 得到，起点为 $I$，
当 $t \geq 0$ 时沿逆时针方向，当 $t < 0$ 时沿顺时针方向。由于
圆周长为 $2\pi$，数 $t$ 与 $t + 2k\pi$（$k \in \Z$）
对应同一点，即 $M(t + 2k\pi) = M(t)$。
\end{definition}

\begin{example}
$M(0) = I$；$M\!\left(\frac{\pi}{2}\right) = (0, 1)$，即圆的最高点；
$M(\pi) = (-1, 0)$；$M\!\left(-\frac{\pi}{2}\right) = (0, -1)$；
且 $M\!\left(\frac{9\pi}{4}\right) = M\!\left(\frac{\pi}{4} +
2\pi\right) = M\!\left(\frac{\pi}{4}\right)$。
\end{example}

\section{余弦与正弦}

\begin{definition}[余弦与正弦]\label{def:g11:trigo:cossin}
对每个\omterm{def:g10:numbers:sets}{实数} $t$，$t$ 的\emph{余弦}\index{余弦}与
\emph{正弦}\index{正弦}分别是
单位圆上点 $M(t)$ 的坐标：
\[
M(t) = (\cos t,\ \sin t).
\]
\end{definition}

\begin{omfigure}
\begin{tikzpicture}[scale=2.0]
  \draw[->] (-1.3,0) -- (1.4,0) node[right] {$x$};
  \draw[->] (0,-1.3) -- (0,1.4) node[above] {$y$};
  \draw[gray] (0,0) circle (1);
  \coordinate (O) at (0,0);
  \coordinate (M) at (55:1);
  \draw[thick,omDef] (O) -- (M);
  \draw[dashed] (M) -- (M |- O) node[below, font=\small] {$\cos t$};
  \draw[dashed] (M) -- (M -| O) node[left, font=\small] {$\sin t$};
  \draw[->,omThm,thick] (0.35,0) arc (0:55:0.35);
  \node[omThm, font=\small] at (27:0.55) {$t$};
  \fill (M) circle (0.7pt) node[above right] {$M(t)$};
  \fill (1,0) circle (0.7pt) node[below right] {$I$};
\end{tikzpicture}

{\small 单位圆：\omterm{def:g10:numbers:sets}{实数} $t$，自
$I$ 出发逆时针卷绕，落在点 $M(t)$ 上，其坐标\emph{定义}了 $\cos t$
和 $\sin t$。}
\end{omfigure}

\begin{proposition}[基本性质]\label{prop:g11:trigo:first}
对任意 $t \in \R$ 及 $k \in \Z$，有
\[
-1 \leq \cos t \leq 1, \qquad
-1 \leq \sin t \leq 1, \qquad
\cos^2 t + \sin^2 t = 1,
\]
\[
\cos(t + 2k\pi) = \cos t, \qquad \sin(t + 2k\pi) = \sin t .
\]
\end{proposition}

\begin{proof}
$M(t)$ 位于单位圆上，故其坐标介于 $-1$ 与
$1$ 之间，且满足圆的\omterm{def:g10:algebra:equation}{方程} $x^2 + y^2 = 1$，这即是
恒等式 $\cos^2 t + \sin^2 t = 1$。周期性
只是重述 $M(t + 2k\pi) = M(t)$（见\cref{def:g11:trigo:winding}）。
\end{proof}

\begin{example}\label{ex:g11:trigo:identity}
若 $\sin t = \frac35$，且 $t \in \intcc{\frac\pi2}{\pi}$（第二
象限），则 $\cos^2 t = 1 - \frac{9}{25} = \frac{16}{25}$，所以
$\cos t = \pm\frac45$；第二象限中\omterm{def:g10:coordgeom:system}{横坐标}为负，
因此 $\cos t = -\frac45$。
\end{example}

需熟记的特殊值：
\begin{center}
\begin{tabular}{c|ccccc}
$t$ & $0$ & $\dfrac{\pi}{6}$ & $\dfrac{\pi}{4}$ & $\dfrac{\pi}{3}$ & $\dfrac{\pi}{2}$\\[6pt]
\hline\rule{0pt}{14pt}%
$\cos t$ & $1$ & $\dfrac{\sqrt3}{2}$ & $\dfrac{\sqrt2}{2}$ & $\dfrac12$ & $0$\\[6pt]
$\sin t$ & $0$ & $\dfrac12$ & $\dfrac{\sqrt2}{2}$ & $\dfrac{\sqrt3}{2}$ & $1$
\end{tabular}
\end{center}

\begin{remark}
记忆口诀：\omterm{def:g11:trigo:cossin}{正弦}值依次为
$\frac{\sqrt0}{2}, \frac{\sqrt1}{2}, \frac{\sqrt2}{2}, \frac{\sqrt3}{2},
\frac{\sqrt4}{2}$，\omterm{def:g11:trigo:cossin}{余弦}值则为该\omterm{def:g11:seq:sequence}{数列}的逆序。
\end{remark}

\section{诱导公式}

\begin{proposition}[诱导公式]\label{prop:g11:trigo:associated}
\index{诱导公式}
对任意 $t \in \R$：
\[
\begin{aligned}
\cos(-t) &= \cos t, & \sin(-t) &= -\sin t,\\
\cos(\pi - t) &= -\cos t, & \sin(\pi - t) &= \sin t,\\
\cos(\pi + t) &= -\cos t, & \sin(\pi + t) &= -\sin t,\\
\cos\left(\tfrac{\pi}{2} - t\right) &= \sin t, &
\sin\left(\tfrac{\pi}{2} - t\right) &= \cos t .
\end{aligned}
\]
\end{proposition}

\begin{proof}
每一行都表示圆的一种对称性。

$M(-t)$ 是 $M(t)$ 关于 $x$ 轴的对称点：\omterm{def:g10:coordgeom:system}{横坐标}
不变，\omterm{def:g10:coordgeom:system}{纵坐标}变号。

$M(\pi - t)$ 是 $M(t)$ 关于 $y$ 轴的对称点：
从 $\pi$ 出发反向走 $t$ 所到之点，与从 $0$ 出发
正向走 $t$ 所到之点水平方向相反。\omterm{def:g10:coordgeom:system}{横坐标}变号，\omterm{def:g10:coordgeom:system}{纵坐标}
不变。

$M(\pi + t)$ 是 $M(t)$ 的对径点（转过半周）：两个
坐标都变号。

$M\!\left(\frac\pi2 - t\right)$ 是 $M(t)$ 关于
直线 $y = x$ 的对称点，交换了两个坐标。
\end{proof}

\begin{omfigure}
\begin{tikzpicture}[scale=2.0]
  \draw[->] (-1.35,0) -- (1.45,0) node[right] {$x$};
  \draw[->] (0,-1.35) -- (0,1.45) node[above] {$y$};
  \draw[gray] (0,0) circle (1);
  \fill (40:1) circle (0.7pt) node[above right, font=\small] {$M(t)$};
  \fill (-40:1) circle (0.7pt) node[below right, font=\small] {$M(-t)$};
  \fill (140:1) circle (0.7pt) node[above left, font=\small]
    {$M(\pi - t)$};
  \fill (220:1) circle (0.7pt) node[below left, font=\small]
    {$M(\pi + t)$};
  \draw[black, dashed, thin] (40:1) -- (-40:1);
  \draw[black, dashed, thin] (40:1) -- (140:1);
  \draw[black, dashed, thin] (140:1) -- (220:1);
  \draw[omDef, thin] (0,0) -- (40:1);
  \draw[omDef, thin] (0,0) -- (140:1);
  \draw[omDef, thin] (0,0) -- (220:1);
  \draw[omDef, thin] (0,0) -- (-40:1);
\end{tikzpicture}

{\small 四个相关点：$M(-t)$ 与 $M(t)$ 关于
$x$ 轴对称，$M(\pi - t)$ 关于 $y$ 轴对称，$M(\pi + t)$
则与之互为对径点。\cref{prop:g11:trigo:associated}中的每一恒等式，
均可由此图读出。}
\end{omfigure}

\begin{example}
$\cos\frac{2\pi}{3} = \cos\left(\pi - \frac\pi3\right)
= -\cos\frac\pi3 = -\frac12$，且
$\sin\frac{7\pi}{6} = \sin\left(\pi + \frac\pi6\right)
= -\sin\frac\pi6 = -\frac12$。五个值的表格经
对称性推广后，便覆盖整个圆周。
\end{example}

\section{解三角方程}

\begin{theorem}[方程 $\cos t = \cos a$ 与 $\sin t = \sin a$]
\label{thm:g11:trigo:equations}
对\omterm{def:g10:numbers:sets}{实数} $t$ 与 $a$：
\[
\cos t = \cos a \iff t = a + 2k\pi \ \text{或}\ t = -a + 2k\pi
\quad (k \in \Z),
\]
\[
\sin t = \sin a \iff t = a + 2k\pi \ \text{或}\ t = \pi - a + 2k\pi
\quad (k \in \Z).
\]
\end{theorem}

\begin{proof}
单位圆上两点的\omterm{def:g10:coordgeom:system}{横坐标}相同，当且仅当它们
相等，或关于 $x$ 轴互为对称点。由
\cref{prop:g11:trigo:associated}可知，$M(a)$ 的对称点为 $M(-a)$。于是
$\cos t = \cos a$ 意味着 $M(t) = M(a)$ 或 $M(t) = M(-a)$，即\
$t = \pm a$（差 $2\pi$ 的\omterm{def:g10:numbers:sets}{整数}倍）。类似地，\omterm{def:g10:coordgeom:system}{纵坐标}相等时，
两点相同，或关于 $y$ 轴互为对称点，而
$M(a)$ 的对称点为 $M(\pi - a)$。
\end{proof}

\begin{method}[在区间上解 $\cos t = c$]\label{met:g11:trigo:solve}
\begin{enumerate}
  \item 由已知值表，求出\emph{一个}角 $a$，
        使 $\cos a = c$。
  \item 写出通解 $t = \pm a + 2k\pi$
        （见\cref{thm:g11:trigo:equations}）。
  \item 选取使解落在所给\omterm{def:g10:numbers:interval}{区间}内的\omterm{def:g10:numbers:sets}{整数} $k$，
        并列出相应的解。
\end{enumerate}
对 $\sin t = c$ 类似处理，其中 $t = a + 2k\pi$ 或
$t = \pi - a + 2k\pi$。
\end{method}

\begin{example}\label{ex:g11:trigo:solve}
求 $\cos t = \frac12$ 在 $\intoc{-\pi}{\pi}$ 上的解。$\cos a = \frac12$ 的一个解为
$a = \frac{\pi}{3}$。通解为
$t = \frac\pi3 + 2k\pi$ 和 $t = -\frac\pi3 + 2k\pi$。在
$\intoc{-\pi}{\pi}$ 内，只有 $k = 0$ 时的解落在其中：
$t \in \left\{-\frac\pi3,\ \frac\pi3\right\}$。

求 $\sin t = -\frac{\sqrt2}{2}$ 在 $\intco{0}{2\pi}$ 上的解。一个角为
$a = -\frac\pi4$，解为 $t = -\frac\pi4 + 2k\pi$ 和
$t = \pi + \frac\pi4 + 2k\pi = \frac{5\pi}{4} + 2k\pi$。在
$\intco{0}{2\pi}$ 内：$t = \frac{7\pi}{4}$（由 $k = 1$ 得到）和
$t = \frac{5\pi}{4}$。
\end{example}

\begin{omfigure}
\begin{tikzpicture}[scale=2.0]
  \draw[->] (-1.35,0) -- (1.45,0) node[right] {$x$};
  \draw[->] (0,-1.2) -- (0,1.3) node[above] {$y$};
  \draw[gray] (0,0) circle (1);
  \draw[omProp, thick] (0.5,-1.05) -- (0.5,1.15)
    node[above, font=\small] {$x = \tfrac12$};
  \fill (60:1) circle (0.7pt)
    node[above right, font=\small] {$M\!\left(\tfrac\pi3\right)$};
  \fill (-60:1) circle (0.7pt)
    node[below right, font=\small] {$M\!\left(-\tfrac\pi3\right)$};
  \draw[omDef, thin] (0,0) -- (60:1);
  \draw[omDef, thin] (0,0) -- (-60:1);
\end{tikzpicture}

{\small 求解 $\cos t = \frac12$：竖直线 $x = \frac12$
（橙色）与单位圆交于两点，这两点关于 $x$ 轴对称
——因此得到两族解 $t = \pm\frac\pi3 + 2k\pi$。}
\end{omfigure}

\section{练习}

\begin{exercise}[$\star$]\label{exo:g11:trigo:1}
将下列角度换算为\omterm{def:g11:trigo:radian}{弧度}：$30^\circ$、$45^\circ$、$120^\circ$、$270^\circ$。
将下列\omterm{def:g11:trigo:radian}{弧度}换算为度数：$\frac{\pi}{5}$、$\frac{5\pi}{6}$、$\frac{7\pi}{4}$、
$\frac{2\pi}{9}$。
\end{exercise}

\begin{exercise}[$\star$]\label{exo:g11:trigo:2}
在单位圆上标出点 $M(t)$，其中
\[
t = \frac{2\pi}{3},\quad
t = -\frac{\pi}{4},\quad
t = \frac{17\pi}{6},\quad
t = -\frac{7\pi}{2},
\]
先将每个值化为 $\intoc{-\pi}{\pi}$ 内的值（模 $2\pi$）。
\end{exercise}

\begin{exercise}[$\star$]\label{exo:g11:trigo:3}
利用表格及诱导公式，求下列各式的精确值：
\[
\cos\frac{3\pi}{4}, \quad
\sin\frac{5\pi}{6}, \quad
\cos\left(-\frac{\pi}{3}\right), \quad
\sin\frac{4\pi}{3}, \quad
\cos\frac{11\pi}{6} .
\]
\end{exercise}

\begin{exercise}[$\star$]\label{exo:g11:trigo:4}
已知 $\cos t = \frac{5}{13}$，且 $t \in \intoo{-\frac\pi2}{0}$，求
$\sin t$ 的精确值。
\end{exercise}

\begin{exercise}[$\star\star$]\label{exo:g11:trigo:5}
对任意\omterm{def:g10:numbers:sets}{实数} $x$，化简：
\[
A = \cos(\pi + x) + \cos(-x) + \sin\left(\frac\pi2 - x\right) - \cos x,
\qquad
B = \sin(\pi - x) + \sin(\pi + x) + \sin(-x) .
\]
\end{exercise}

\begin{exercise}[$\star\star$]\label{exo:g11:trigo:6}
在 $\intoc{-\pi}{\pi}$ 上解下列\omterm{def:g10:algebra:equation}{方程}：
\[
\cos t = \frac{\sqrt3}{2}, \qquad
\sin t = \frac12, \qquad
\cos t = -1 .
\]
\end{exercise}

\begin{exercise}[$\star\star$]\label{exo:g11:trigo:7}
在 $\intco{0}{2\pi}$ 上解下列\omterm{def:g10:algebra:equation}{方程}：
\[
\sin t = -\frac{\sqrt3}{2}, \qquad
\cos t = 0, \qquad
2\sin t + 1 = 0 .
\]
\end{exercise}

\begin{exercise}[$\star\star$]\label{exo:g11:trigo:8}
在 $\R$ 上解\omterm{def:g10:algebra:equation}{方程} $\cos 2t = \cos t$。（利用
\cref{thm:g11:trigo:equations}，取 $a = t$，并讨论两族
解的情况。）
\end{exercise}

\begin{exercise}[$\star\star$]\label{exo:g11:trigo:9}
证明：对所有\omterm{def:g10:numbers:sets}{实数} $x$，有
\[
\bigl(\cos x + \sin x\bigr)^2 + \bigl(\cos x - \sin x\bigr)^2 = 2 .
\]
\end{exercise}

\begin{exercise}[$\star\star$]\label{exo:g11:trigo:10}
一个半径为 $30$ cm 的车轮做无滑动滚动。自行车前进 $1.5$ m 时，车轮转过的角度
（分别以\omterm{def:g11:trigo:radian}{弧度}和度为单位）是多少？
车轮转动一整圈时，自行车前进多远？
\end{exercise}

\begin{exercise}[$\star\star\star$]\label{exo:g11:trigo:11}
在\omterm{def:g10:numbers:interval}{区间} $\intoc{-\pi}{\pi}$ 上解不等式
$\cos t \leq \frac12$。（画出单位圆，标出
\omterm{def:g10:coordgeom:system}{横坐标}不超过 $\frac12$ 的区域，并读出
$t$ 的取值\omterm{def:g10:numbers:interval}{区间}。）
\end{exercise}

\section{问题：摩天轮与潮汐}

\begin{problem}[{周末问题——弧度即滚过的距离，
以及单位圆作为一切周期现象的
主时钟}]
\label{pb:g11:trigo:1}
摩天轮载着你绕圆周运行；潮汐把港口的海水
沿着同一个数学圆周涨落。一切
周期现象——车轮、潮汐、声音、季节——都从
本章的单位圆上读出自己的时刻表
（见\cref{def:g11:trigo:cossin}），而
\cref{met:g11:trigo:solve}中的\omterm{def:g10:algebra:equation}{方程}给出准确时刻。本题
将依次进行换算、绕行、建模与求解，并在此过程中
说明数学家为何用\omterm{def:g11:trigo:radian}{弧度}度量角。

\medskip
\noindent\textbf{第一部分——\omterm{def:g11:trigo:radian}{弧度}，滚过的距离。}
\begin{enumerate}
  \item 换算为\omterm{def:g11:trigo:radian}{弧度}：$30^\circ$、$135^\circ$、
        $300^\circ$；换算为度：$\frac{3\pi}{4}$、
        $\frac{5\pi}{6}$（见\cref{met:g11:trigo:convert}）。
  \item \omterm{def:g11:trigo:radian}{弧度}的要点：圆心角为 $t$（以\omterm{def:g11:trigo:radian}{弧度}计）的弧，
        在半径为 $r$ 的圆上，弧长恰好为
        $r\,t$。取圆心角为 $2$ \omterm{def:g11:trigo:radian}{弧度}、半径为 $3$ m 的圆，
        它截出的弧长是多少？（全程不出现 $\pi$——这正是
        要点。）
  \item 半径为 $30$ cm 的车轮滚过 $1$ km（无滑动）。
        它共转过多少\emph{\omterm{def:g11:trigo:radian}{弧度}}，
        合多少整圈？
  \item 在单位圆上标出并精确求值：
        $\cos\frac{2\pi}{3}$、$\sin\left(-\frac{\pi}{4}\right)$、
        以及 $\cos\frac{19\pi}{6}$（先模 $2\pi$ 化简；
        见\cref{prop:g11:trigo:associated}）。
  \item 角 $t \in \intoo{\frac\pi2}{\pi}$ 满足
        $\sin t = \frac35$。利用
        $\cos^2 t + \sin^2 t = 1$ 及象限信息，求
        $\cos t$ 与 $\tan t$。
\end{enumerate}

\medskip
\noindent\textbf{第二部分——摩天轮。}
一座巨型摩天轮半径为 $60$ m，轮轴高出
地面 $65$ m，每 $30$ 分钟转一圈。你在
最低点登上吊舱，此时 $t = 0$（单位：分钟），因此你的高度满足
\[
h(t) = 65 - 60 \cos\!\left(\frac{\pi t}{15}\right).
\]
\begin{enumerate}[resume]
  \item 说明该公式：验证经过 $t$
        分钟后转过的角度、$t = 0$ 时的高度，并解释
        负号的含义。
  \item 求你在 $t = 7.5$、$t = 15$ 及
        $t = 20$ 分钟时的高度（精确值）。
  \item 你于何时\emph{首次}到达 $95$ m 高？（利用
        \cref{met:g11:trigo:solve}解\omterm{def:g10:algebra:equation}{方程} $h(t) = 95$。）
  \item 在乘坐过程的哪一时间\omterm{def:g10:numbers:interval}{区间}内，你的高度至少为
        $95$ m——总共持续多长时间？（可与
        \cref{exo:g11:trigo:11}比较。）
  \item 高度在何处变化\emph{最快}？比较
        第一分钟内的上升量
        （$h(1) - h(0)$）与第 $7$ 至
        $8$ 分钟之间的上升量，并归纳这一观察（使之精确的
        工具——\omterm{def:g11:trigo:cossin}{正弦}型\omterm{def:g11:func:function}{函数}的\omterm{def:g11:deriv:derivative}{导数}——
        将于十二年级学到）。
  \item 工程挑战：设计一座摩天轮的高度公式，
        其最大高度为 $40$ m，最小高度为 $4$ m，
        周期为 $12$ 分钟，乘客在
        $t = 0$ 时于最低点登上吊舱。
\end{enumerate}

\medskip
\noindent\textbf{第三部分——潮汐与圆的
对称性。}
在某港口，午夜后 $t$ 小时的水深（以米计）
为
\[
d(t) = 7 + 3 \sin\!\left(\frac{\pi t}{6}\right).
\]
\begin{enumerate}[resume]
  \item 给出最大水深与最小水深，以及它们首次出现的时刻，
        并求潮汐的周期。
  \item 一艘货船需要 $8.5$ m 的水深。求
        不等式 $d(t) \geq 8.5$ 在 $\intco{0}{12}$ 上的解：货船
        可在哪一时段进港？该时段有多长？下一个
        可进港时段何时开启？
  \item 直接由圆的对称性重新推导两个诱导公式——
        $\sin(\pi - t) = \sin t$（镜面对称，
        沿纵轴）与 $\cos(\pi + t) = -\cos t$
        （半圈旋转）——并写出余角恒等式
        $\cos\left(\frac\pi2 - t\right) = \sin t$。
  \item 奇偶性（术语见\cref{pb:g11:func:1}）：判断
        $\sin$、$\cos$ 与 $t \mapsto \sin t + t^3$ 是\omterm{def:g11:func:parity}{偶函数}还是\omterm{def:g11:func:parity}{奇函数}，
        并用单位圆各用一句话说明理由。
  \item 在 $\intco{0}{2\pi}$ 上解\omterm{def:g10:algebra:equation}{方程}：
        $2\sin^2 t - \sin t - 1 = 0$。（伪装的
        二次\omterm{def:g10:algebra:equation}{方程}：按\cref{pb:g11:quad:1}的方法\omterm{def:g10:algebra:expand}{因式分解}，再
        读单位圆上的解。）
\end{enumerate}

\medskip
\noindent\textbf{第四部分——数学家的角度。}
\begin{enumerate}[resume]
  \item $1$ \omterm{def:g11:trigo:radian}{弧度}等于多少度？再看一个计算器陷阱：
        $\sin(1)$（\omterm{def:g11:trigo:radian}{弧度}模式）与
        $\sin(1^\circ)$ 相比哪个更大，大约相差多少倍？这对
        计算器的设置有何启示？
  \item 求 $\cos t = \sin t$ 在 $\intco{0}{2\pi}$ 上的解，
        利用单位圆（\omterm{def:g10:coordgeom:system}{横坐标}与\omterm{def:g10:coordgeom:system}{纵坐标}在何处相等？）。
  \item 对小角而言，弧与其\omterm{def:g11:trigo:cossin}{正弦}值几乎重合：
        比较 $\sin(0.1)$ 与 $0.1$（保留五位小数），并
        给出相对误差。这种“小角
        近似”用于引导船舶与单摆；其背后的定理
        （$\frac{\sin t}{t} \to 1$）是
        十二年级极限一章的重点定理。
  \item 总结——单位圆作为主时钟，每项一句话：
        \omterm{def:g11:trigo:radian}{弧度}把角化为确定的数（弧长
        $=$ 半径 $\times$ 角）；\omterm{def:g11:trigo:cossin}{余弦}与\omterm{def:g11:trigo:cossin}{正弦}分别是
        时针的影子与高度；圆的对称性
        即恒等式；解 $\cos t = c$ 即从时钟上读出
        时刻；而一切周期现象——车轮、
        潮汐、声音——都不过是穿了外衣的同一只时钟。
\end{enumerate}
\end{problem}
